% Generated draft: requires mathematical and source review.
\documentclass[11pt]{article}
\usepackage[margin=25mm]{geometry}
\usepackage{fontspec}
\setmainfont{Latin Modern Roman}
\usepackage{amsmath,amssymb,mathtools,unicode-math}
\setmathfont{Latin Modern Math}
\usepackage{enumitem,needspace,xurl,hyperref}
\setlength{\parindent}{0pt}
\setlength{\parskip}{5pt}
\setlength{\emergencystretch}{3em}
\newcommand{\N}{\mathbb{N}}
\newcommand{\Z}{\mathbb{Z}}
\newcommand{\Q}{\mathbb{Q}}
\newcommand{\R}{\mathbb{R}}
\newcommand{\C}{\mathbb{C}}
\newcommand{\F}{\mathbb{F}}
\newcommand{\A}{\mathbb{A}}
\newcommand{\PP}{\mathbb{P}}
\newcommand{\E}{\mathbb{E}}
\newcommand{\eps}{\varepsilon}
\newcommand{\dd}{\,\mathrm{d}}
\begin{document}
\section*{UnsolvedMath ID: 700899}
\textbf{Source problem code:} catalogue entry
\subsection{Definitions and mathematical statement}
Let \(p\) be a prime. A variety of groups is a class of groups defined by group identities. A subvariety of a variety \(\mathcal V\) is a variety \(\mathcal W\) satisfying \(\mathcal W\subseteq\mathcal V\). Write
\[
\mathfrak A_p=\{G:G\text{ is abelian and }g^p=1\text{ for every }g\in G\}.
\]
Thus the exponent divides \(p\); the trivial group is included, as required for a variety. No finiteness assumption is imposed.

For varieties \(\mathcal U\) and \(\mathcal V\), their product \(\mathcal U\mathcal V\) consists of the groups \(G\) having a normal subgroup \(N\trianglelefteq G\) with \(N\in\mathcal U\) and \(G/N\in\mathcal V\). Under the usual interpretation of the superscript as a product of varieties, put
\[
\mathfrak A_p^3=(\mathfrak A_p\mathfrak A_p)\mathfrak A_p,
\qquad
\operatorname{Sub}(\mathfrak A_p^3)
=\{\mathcal W:\mathcal W\text{ is a variety of groups and }\mathcal W\subseteq\mathfrak A_p^3\}.
\]
\textbf{Problem.} For each prime \(p\), determine the cardinality
\[
\bigl|\operatorname{Sub}(\mathfrak A_p^3)\bigr|.
\]
The supplied statement does not explicitly define the superscript \(3\). The formulation above uses the threefold product convention, which requires confirmation against the original source. The source's phrase ``prime exponent \(p\)'' is interpreted by the identity \(x^p=1\), rather than by excluding the trivial group.
\subsection{Short English statement}
For a prime \(p\), determine how many distinct subvarieties the threefold product of the variety of abelian groups of exponent dividing \(p\) has. The product interpretation of \(\mathfrak A_p^3\) requires source confirmation. The supplied registry labels the problem open; this draft does not independently verify that status.
\subsection{Sources}
\begin{itemize}
\item Local UnsolvedMath registry, supplied complete source record: permanent ID \(700899\), exact source problem code \texttt{catalogue entry}, title \emph{Kourovka Notebook Problem 12.81}. The record was created and updated at \texttt{2026-10-16T14:00:00Z} and carries the status \texttt{open}; these are supplied registry data.
\item \emph{The Kourovka Notebook (Unsolved Problems in Group Theory)}, problem 12.81, 12th issue (1992); the supplied background also cites the 21st issue edition (2026), \texttt{arXiv:1401.0300v48}. The collection description names E. I. Khukhro and V. D. Mazurov as editors. These bibliographic details are supplied attribution and have not been independently checked.
\item Proposed in 1992. The proposer field is malformed; its exact supplied text is \texttt{V. I. Sushchanski\textbackslash{}u\{\textbackslash{}i}. No repair of the name is asserted.
\end{itemize}
\end{document}
